% Le petit dictionnaire : chaque mot en langage simple.
\chapter{Le petit dictionnaire}
\markboth{Le petit dictionnaire}{Le petit dictionnaire}

Chaque mot est expliqué en une ou deux phrases, sans jargon. Les mots en italique dans une définition ont leur propre
entrée.

\begin{description}[style=nextline,leftmargin=0pt,font=\normalfont\sffamily\bfseries\color{okbleu},itemsep=5pt]
\item[Agent d'IA] Un programme d'intelligence artificielle à qui l'on confie une tâche entière (chercher, écrire du
code, calculer, rédiger) et qui l'exécute seul, étape par étape. Il peut se tromper en disant qu'il a réussi.
\item[Arithmétique de boules (ou d'intervalles)] Une façon de calculer où chaque nombre est remplacé par un petit
intervalle dont on est \emph{sûr} qu'il contient la vraie valeur. Plus lente que le calcul ordinaire, elle donne des
garanties au lieu d'approximations.
\item[Audit] L'examen d'un travail par un agent qui ne l'a pas produit : il recompile, refait les calculs et rend un
verdict écrit. Dans ce livre, les auditeurs sont des agents d'IA ; un audit n'est donc pas une \emph{relecture par des
pairs}.
\item[Axiome] Une règle que l'on admet sans la démontrer. Lean accepte trois règles de base, toujours les mêmes ; en
ajouter une autre en douce permet de « démontrer » n'importe quoi, même $1=2$.
\item[BAO (oscillations acoustiques des baryons)] Une « règle graduée » laissée dans la répartition des galaxies par
les ondes sonores de l'univers très jeune. Elle permet de mesurer de très grandes distances.
\item[Certificat] Un ensemble de données qui permet de vérifier rapidement un résultat difficile à trouver. Comme une
solution de grille de sudoku : longue à chercher, rapide à contrôler.
\item[Conjecture] Une affirmation que l'on croit vraie, mais que personne n'a encore démontrée.
\item[Contrôle (positif, négatif)] Un cas test dont on connaît la réponse à l'avance. Le contrôle positif doit
répondre « oui », le négatif « non ». S'ils ne répondent pas ainsi, c'est le test qui est cassé.
\item[Densité] Le nombre de points par unité de surface. Un réseau « de densité~1 » a, en moyenne, un point par carré
de côté~1.
\item[DOI] Un identifiant permanent attribué à un travail publié, comme un numéro d'ISBN pour un livre. Il permet de
le retrouver même si son adresse change.
\item[Énergie de point zéro] Même au zéro absolu, dans le monde quantique, les particules ne sont jamais parfaitement
immobiles ; elles vibrent un peu. L'énergie de cette vibration minimale s'appelle l'énergie de point zéro.
\item[Fitness] En biologie de l'évolution, une mesure du succès reproductif d'un individu ou d'un type : combien de
descendants il laisse, en moyenne.
\item[Fixation] On dit qu'une mutation se « fixe » quand elle finit par être portée par toute la population. La
probabilité de fixation est la chance que cela arrive.
\item[Glose] La phrase en français (ou en anglais) qui dit ce qu'un énoncé Lean est censé vouloir dire. Le juge
vérifie l'énoncé, pas la glose.
\item[Graphe isotherme] Un réseau d'individus (qui peut remplacer qui) où chacun est remplacé, au total, au même
rythme. Le théorème isotherme dit qu'un tel réseau ne favorise ni ne gêne la sélection naturelle.
\item[Lean] Un langage informatique dans lequel on écrit des démonstrations mathématiques assez détaillées pour qu'un
ordinateur puisse les vérifier.
\item[Neuro-symbolique] Une façon de faire travailler ensemble des réseaux de neurones, qui proposent, et des méthodes
logiques, qui vérifient.
\item[Niveau (A, L, C, X)] L'étiquette de confiance d'une affirmation dans ce livre : A, « le juge a regardé » ; L,
« c'est écrit dans un livre sérieux » ; C, « un raisonnement, mené ici par des agents d'IA, que personne n'a vérifié à la machine » ; X, « une
esquisse au crayon ».
\item[Noyau] Le petit programme au centre de Lean qui vérifie chaque pas d'une démonstration. C'est le « juge » de ce
livre.
\item[Nombre à virgule (flottante)] La façon ordinaire dont un ordinateur représente les nombres réels, avec un
nombre limité de chiffres. Chaque opération arrondit un peu ; d'où le niveau X.
\item[openai/math] Une publication d'OpenAI contenant des démonstrations produites par une IA et vérifiées par
ordinateur, dont le théorème sur le réseau triangulaire. Elle n'a pas été relue par des pairs.
\item[Potentiel (ou force) répulsif] La règle qui dit combien deux particules se repoussent selon leur distance.
Certaines familles de potentiels, dites « complètement monotones », sont couvertes par le théorème sur le réseau
triangulaire.
\item[Pré-enregistrement] Écrire et archiver, avant une expérience, ce que l'on prédit et ce qui prouverait qu'on a
tort.
\item[Preuve (démonstration)] Un enchaînement de pas logiques qui mène, sans trou, des hypothèses à la conclusion.
\item[Registre des affirmations] Un fichier qui contient une ligne par affirmation d'un article, avec son niveau, sa
preuve et la case « audit humain ». Un programme le lit et dit « non » tant qu'il manque quelque chose.
\item[Rejeu] La relecture complète, par un second programme (\texttt{leanchecker}), de toutes les démonstrations
enregistrées, en les soumettant de nouveau au noyau. Il attrape les pièces glissées sans contrôle.
\item[Relecture par des pairs] L'examen d'un travail par des chercheurs humains du domaine, avant ou après sa
publication. Aucun résultat de ce livre ne l'a encore reçue.
\item[Réseau (cristallin)] Une disposition de points qui se répète à l'identique dans toutes les directions, comme les
alvéoles d'une ruche ou les carreaux d'un carrelage.
\item[Réseau triangulaire] Le réseau où chaque point a six voisins à égale distance, disposés en triangles. Ses
cellules de Voronoï sont des hexagones, comme sur la couverture.
\item[Résultat négatif] Le résultat d'un test qui montre qu'une idée ne marche pas. Il est publié dans ce livre avec
le même soin qu'un succès.
\item[RK4] Une méthode très répandue pour faire avancer une simulation pas à pas dans le temps.
\item[Simulation] Un calcul par ordinateur qui imite l'évolution d'un système (des particules, un fluide, une
population) en suivant ses équations.
\item[Solveur raide] Un programme qui résout des équations où certaines choses changent très vite et d'autres très
lentement, comme en chimie ou dans un plasma. Les méthodes ordinaires y deviennent instables ; un solveur raide, comme
CVODE et sa traduction rusty-SUNDIALS, reste stable.
\item[sorry] Le mot qu'on écrit en Lean à la place d'un morceau de démonstration qu'on n'a pas fini. Le fichier se
compile quand même : c'est pourquoi « compilé sans erreur » ne veut pas dire « démontré ».
\item[Stellarator] Une machine en forme d'anneau tordu, conçue pour confiner un gaz très chaud par des champs
magnétiques, en vue de la fusion nucléaire.
\item[Superfluide] Un liquide qui, à très basse température, coule sans aucune friction, comme l'hélium
refroidi tout près du zéro absolu.
\item[Tension de Hubble] Le désaccord entre deux familles de mesures de la vitesse d'expansion de l'univers.
\item[Voronoï (cellule de)] Pour un point d'un ensemble, la région du plan plus proche de ce point que de tous les
autres.
\item[Yukawa (force de)] Une force répulsive « écrantée » : elle s'affaiblit beaucoup plus vite que la force
électrique ordinaire, parce que le milieu environnant la masque. On la rencontre dans les plasmas poussiéreux et les
suspensions colloïdales.
\item[Zenodo] Une archive scientifique ouverte où chacun peut déposer un travail et recevoir un DOI.
\end{description}
